Clause (a): the current's positivity, conservation and uniqueness are Lemma \claimref{one-relation.dirac-current-conserved-density}.

Clause (b). The blade inverse is the anti-automorphism that fixes the generators of square $1$ and negates those of square $-1$ (Premise \claimref{synthesis.premise-positivity-of-values}); $x \mapsto \gamma_{0}\tilde{x}\gamma_{0}$ is an anti-automorphism, fixes $\gamma_{0}$ and sends $\gamma_{i}$ to $\gamma_{0}\gamma_{i}\gamma_{0} = -\gamma_{i}$, so the two agree. Hence $\langle\psi^{\dagger}\psi\rangle_{0} = \langle\gamma_{0}\tilde{\psi}\gamma_{0}\psi\rangle_{0} = \langle\psi\gamma_{0}\tilde{\psi}\gamma_{0}\rangle_{0} = J \cdot \gamma_{0}$ by cyclicity of the scalar part, $J$ being a vector for even $\psi$. The blades of square $1$ other than $1$ are $\gamma_{0}$, the three $\gamma_{0}\gamma_{i}$ and $\gamma_{1}\gamma_{2}\gamma_{3}$; the last four anticommute with $\gamma_{0}$, so $\{\gamma_{0}\}$ is a maximal independent commuting set and $\Pi = \tfrac12(1 + \gamma_{0})$ is an atom by clause (i) of Theorem \claimref{synthesis.born-weights-under-positivity-a-state-sharp-on-an-atom-is-its-trace-functional}. With $\Pi = \Pi^{\dagger} = \Pi^{2}$, $\langle(\psi\Pi)^{\dagger}(\psi\Pi)\rangle_{0} = \langle\Pi\psi^{\dagger}\psi\Pi\rangle_{0} = \langle\psi^{\dagger}\psi\Pi\rangle_{0} = \tfrac12\langle\psi^{\dagger}\psi\rangle_{0} + \tfrac12\langle\psi^{\dagger}\psi\gamma_{0}\rangle_{0}$, and the last term is the scalar part of an odd element, zero. The script confirms the three identities and the maximality on four sample spinors.

Clause (c). By clause (ii) of the theorem just cited $\Pi_{\psi}$ is an atom with $n_{\psi}\langle\Pi\rangle_{0} = \langle\Pi\psi^{\dagger}\psi\rangle_{0} = \langle(\psi\Pi)^{\dagger}(\psi\Pi)\rangle_{0}$, and by its clause (iii), which is where (P) is used, a state sharp on $\Pi_{\psi}$ is $x \mapsto \langle\Pi_{\psi}x\rangle_{0}/\langle\Pi\rangle_{0}$. For a projector $E$, cyclicity and $E = E^{\dagger}E$, $\Pi = \Pi^{2}$ give $\langle\psi\Pi\psi^{\dagger}E\rangle_{0} = \langle\Pi\psi^{\dagger}E^{\dagger}E\psi\Pi\rangle_{0} = \langle(E\psi\Pi)^{\dagger}(E\psi\Pi)\rangle_{0}$, and the displayed formula follows on dividing. For orthogonal projectors with sum $1$ the weights sum to $\omega(1) = 1$. The instances are computed by the script from both sides of the formula. That a state sharp on the atom is this functional is the cited theorem and is not recomputed.

Clause (d). The indifference step is Fact \claimref{one-relation.imported-indifference-instrument} applied to components related by an exact symmetry of the vacuum, which is precisely the instrument's licensed use (equating along orbits), and \emph{not} the illicit use audited in the families paper (equating inequivalent components). The script exhibits the licensed step in miniature: on $n = 2, 3, 4, 5$ equal-norm components the weight vectors invariant under all transpositions form a one-dimensional space, the uniform one, while under an intransitive subgroup (one transposition, or two disjoint ones, on four components) the invariant space has dimension three or two, one per orbit, and under the transitive Klein four-group it is again one-dimensional: equality is forced by transitivity on the components, which is what an exact orbit supplies.

The residual postulate. That (P) is equivalent to non-negative, normalised, additive values on effects is clause (ii) of Theorem \claimref{synthesis.positivity-of-values-is-additivity-of-non-negative-weights-on-effects}; the identification of that statement with the residual postulate of versions 1 to 3 is the reading argued there and is stated, not proved.

Prior art. The fine-graining and continuity steps of versions 1 to 3 were the envariance argument's \cite{Zurek2005}; they are not steps of this proof. The script still computes the fine-graining count on three rational instances, as a comparison: given additivity it returns the squared norms that clause (c) gives.

Scope item (4). For $r = \gamma_{2}\gamma_{3}$: $r$ commutes with $\gamma_{0}$, hence with $\Pi$, and $rr^{\dagger} = 1$, so $(\psi r)\Pi(\psi r)^{\dagger} = \psi\Pi rr^{\dagger}\psi^{\dagger} = \psi\Pi\psi^{\dagger}$ and $n_{\psi r} = n_{\psi}$; the script confirms it and that the vectors $\psi\gamma_{3}\tilde{\psi}$ and $\psi r\gamma_{3}\tilde{r}\tilde{\psi}$ differ on the four samples. That spin projectors act through right multiplication is the usual description of spin for an even spinor and is argued, not proved here.
